% GENERATED FILE. Edit canonical lessons, then run npm run generate:books.
% canonical-source-hash: fnv1a64:4803ac7c9116060e
% canonical-lessons: IT-C473-sincero, IT-C473-bugiardo, IT-C473-ingenuo, IT-C473-furbo, IT-C473-ammirare

\chapter{Sinceri e bugiardi}
\label{ch:it-sinceri-e-bugiardi}
\hlchaptermodality{473}

\begin{chapteropening}
\textbf{By the end of this chapter:} \emph{I can say whether someone is honest, naive or sly.}

Complete the last lesson of chapter 473: I can say whether someone is honest, naive or sly.
\end{chapteropening}

\section[sincero]{sincero — sincere, honest}

\label{lesson:IT-C473-sincero}



\begin{quote}
Before the new one: say the Italian for mean, then the Italian for to boast.
\end{quote}

\subsection*{You'll want to know: sincero}
\textbf{sincero} — \textquotedblleft{}sincere, honest\textquotedblright{}.

\begin{quote}
Sii sincero: ti piace il mio vestito?
\end{quote}

\begin{quote}
\emph{Be honest: do you like my dress?}
\end{quote}

The noun is \emph{la sincerità}.

\begin{cousinweb}
\emph{sincero} is a learned borrowing from Latin \emph{sincērus} \textquotedblleft{}pure, clean, whole\textquotedblright{} (EtymDB; Lewis \& Short).
\end{cousinweb}

\begin{grammarlens}[title={What you've built}]
One more word for telling, explaining and describing at B1.
\end{grammarlens}

\subsection*{Your turn}
\begin{itemize}
\raggedright
  \item \emph{Say it:} \emph{sincero}
  \item \emph{Say it:} \emph{sincero}, once more
  \item \emph{Say it:} a sentence of your own with \emph{sincero}
  \item \emph{Recall:} say the Italian for mean, then the Italian for to boast, then say \emph{sincero} again
\end{itemize}

\begin{tcolorbox}[breakable,colback=teal!4,colframe=teal!35!black,title={Before you move on}]
What does sincero mean? (\textquotedblleft{}Sincere, honest\textquotedblright{}.) Say it once more.
\end{tcolorbox}
\section[il bugiardo]{il bugiardo — liar; lying}

\label{lesson:IT-C473-bugiardo}



\begin{quote}
Before the new one: say the Italian for to boast, then the Italian for sincere.
\end{quote}

\subsection*{You'll want to know: il bugiardo}
\textbf{il bugiardo} — \textquotedblleft{}liar; lying\textquotedblright{}.

\begin{quote}
Non credergli, è un bugiardo.
\end{quote}

\begin{quote}
\emph{Do not believe him, he is a liar.}
\end{quote}

Built on \emph{la bugia} \textquotedblleft{}a lie\textquotedblright{}. \emph{Dire bugie} — to tell lies.

\begin{cousinweb}
\emph{bugiardo} is built on \emph{bugia} \textquotedblleft{}a lie\textquotedblright{} with \emph{-ardo} (EtymDB). \emph{Bugia} is usually traced to Old Occitan \emph{bauzia} \textquotedblleft{}deceit\textquotedblright{}, itself from a Germanic word.
\end{cousinweb}

\begin{grammarlens}[title={What you've built}]
One more word for telling, explaining and describing at B1.
\end{grammarlens}

\subsection*{Your turn}
\begin{itemize}
\raggedright
  \item \emph{Say it:} \emph{il bugiardo}
  \item \emph{Say it:} \emph{il bugiardo}, once more
  \item \emph{Say it:} a sentence of your own with \emph{il bugiardo}
  \item \emph{Recall:} say the Italian for to boast, then the Italian for sincere, then say \emph{il bugiardo} again
\end{itemize}

\begin{tcolorbox}[breakable,colback=teal!4,colframe=teal!35!black,title={Before you move on}]
What does il bugiardo mean? (\textquotedblleft{}Liar; lying\textquotedblright{}.) Say it once more.
\end{tcolorbox}
\section[ingenuo]{ingenuo — naive}

\label{lesson:IT-C473-ingenuo}



\begin{quote}
Before the new one: say the Italian for sincere, then the Italian for liar.
\end{quote}

\subsection*{You'll want to know: ingenuo}
\textbf{ingenuo} — \textquotedblleft{}naive\textquotedblright{}.

\begin{quote}
Ero così ingenuo che credevo a tutto.
\end{quote}

\begin{quote}
\emph{I was so naive that I believed everything.}
\end{quote}

Do not mix it up with \emph{ingegnoso} \textquotedblleft{}clever, ingenious\textquotedblright{}: \emph{ingenuo} means naive, like English \emph{ingenuous}.

\begin{cousinweb}
\emph{ingenuo} is a learned borrowing from Latin \emph{ingenuus} \textquotedblleft{}native, freeborn\textquotedblright{}, and so \textquotedblleft{}frank, open\textquotedblright{} (Lewis \& Short).
\end{cousinweb}

\begin{grammarlens}[title={What you've built}]
One more word for telling, explaining and describing at B1.
\end{grammarlens}

\subsection*{Your turn}
\begin{itemize}
\raggedright
  \item \emph{Say it:} \emph{ingenuo}
  \item \emph{Say it:} \emph{ingenuo}, once more
  \item \emph{Say it:} a sentence of your own with \emph{ingenuo}
  \item \emph{Recall:} say the Italian for sincere, then the Italian for liar, then say \emph{ingenuo} again
\end{itemize}

\begin{tcolorbox}[breakable,colback=teal!4,colframe=teal!35!black,title={Before you move on}]
What does ingenuo mean? (\textquotedblleft{}Naive\textquotedblright{}.) Say it once more.
\end{tcolorbox}
\section[furbo]{furbo — sly, crafty; clever}

\label{lesson:IT-C473-furbo}



\begin{quote}
Before the new one: say the Italian for liar, then the Italian for naive.
\end{quote}

\subsection*{You'll want to know: furbo}
\textbf{furbo} — \textquotedblleft{}sly, crafty; clever\textquotedblright{}.

\begin{quote}
Il gatto è furbo: apre la porta da solo.
\end{quote}

\begin{quote}
\emph{The cat is crafty: it opens the door by itself.}
\end{quote}

Often a compliment in speech: \emph{che furbo!} — how clever!

\begin{cousinweb}
\emph{furbo} was borrowed from French \emph{fourbe} \textquotedblleft{}a cheat\textquotedblright{} (EtymDB), which is built on \emph{fourbir} \textquotedblleft{}to polish\textquotedblright{}, a verb of Germanic origin.
\end{cousinweb}

\begin{grammarlens}[title={What you've built}]
One more word for telling, explaining and describing at B1.
\end{grammarlens}

\subsection*{Your turn}
\begin{itemize}
\raggedright
  \item \emph{Say it:} \emph{furbo}
  \item \emph{Say it:} \emph{furbo}, once more
  \item \emph{Say it:} a sentence of your own with \emph{furbo}
  \item \emph{Recall:} say the Italian for liar, then the Italian for naive, then say \emph{furbo} again
\end{itemize}

\begin{tcolorbox}[breakable,colback=teal!4,colframe=teal!35!black,title={Before you move on}]
What does furbo mean? (\textquotedblleft{}Sly, crafty; clever\textquotedblright{}.) Say it once more.
\end{tcolorbox}
\section[ammirare]{ammirare — to admire}

\label{lesson:IT-C473-ammirare}



\begin{quote}
Before the new one: say the Italian for naive, then the Italian for sly.
\end{quote}

\subsection*{You'll want to know: ammirare}
\textbf{ammirare} — \textquotedblleft{}to admire\textquotedblright{}.

\begin{quote}
Ammiro molto il coraggio di mia madre.
\end{quote}

\begin{quote}
\emph{I really admire my mother's courage.}
\end{quote}

\emph{Ammirare il panorama} — to take in the view.

\begin{cousinweb}
\emph{ammirare} is a learned borrowing from Latin \emph{admīrārī} \textquotedblleft{}to wonder at\textquotedblright{} (EtymDB), built on \emph{mīrārī} \textquotedblleft{}to wonder\textquotedblright{} (Lewis \& Short).
\end{cousinweb}

\begin{grammarlens}[title={What you've built}]
One more word for telling, explaining and describing at B1.
\end{grammarlens}

\subsection*{Your turn}
\begin{itemize}
\raggedright
  \item \emph{Say it:} \emph{ammirare}
  \item \emph{Say it:} \emph{ammirare}, once more
  \item \emph{Say it:} a sentence of your own with \emph{ammirare}
  \item \emph{Recall:} say the Italian for naive, then the Italian for sly, then say \emph{ammirare} again
\end{itemize}

\begin{tcolorbox}[breakable,colback=teal!4,colframe=teal!35!black,title={Before you move on}]
What does ammirare mean? (\textquotedblleft{}To admire\textquotedblright{}.) Say it once more.
\end{tcolorbox}
